\section{Cơ sở nghiên cứu và định vị kiến trúc}
\label{sec:background}

\subsection{Kiến trúc tham chiếu, hệ sinh thái nhiều chủ thể và truy vết}
\label{subsec:ra}

Kiến trúc tham chiếu phần mềm (SRA) mô tả kiến trúc tổng quát cho một lớp hệ thống, thay vì một cấu hình triển khai cụ thể \citep{angelov2012}. OASIS SOA-RAF cũng xem kiến trúc tham chiếu là tập phần tử trừu tượng của một miền, độc lập với sản phẩm và đặc biệt quan tâm tới tương tác qua ranh giới sở hữu \citep[pp.~9--14]{oasis_soaraf2012}. Theo ISO/IEC/IEEE 42010:2022, mô tả kiến trúc được tổ chức theo bên liên quan, mối quan tâm, góc nhìn và mô hình \citep{iso42010_2022}. Vì vậy, nghiên cứu này coi SRA là tập các quyết định và bất biến cần được bảo toàn qua nhiều hiện thực hóa, còn topology lưu trữ, cơ chế tích hợp và công nghệ thực thi là các lựa chọn có thể biến thiên.

Cách tiếp cận này đặc biệt phù hợp với hệ sinh thái phần mềm nhiều chủ thể. SRA có thể đóng vai trò điều phối cả quan hệ kỹ thuật và tổ chức giữa các thành phần tự chủ \citep{knodel2016}; các nghiên cứu về hệ sinh thái nhiều nhà cung cấp cũng sử dụng kiến trúc tham chiếu để làm rõ vai trò, giao diện và cách ánh xạ các cấu hình thực tế vào một cấu trúc chung \citep{kruize2016}. Trong bài toán tỉnh--trung ương, các tỉnh, nền tảng quốc gia và CSDL chuyên ngành có chủ thể vận hành khác nhau, nên kiến trúc cần mô tả rõ trách nhiệm và thẩm quyền mà không đồng nhất chúng với vị trí vật lý của dữ liệu.

Truy vết được dùng để nối căn cứ chuẩn tắc với quyết định kiến trúc. Breaux và cộng sự cho thấy quyền và nghĩa vụ trong quy định cần được trích xuất và căn chỉnh với yêu cầu hệ thống \citep{breaux2006}; các nghiên cứu gần đây tiếp tục nhấn mạnh việc đưa tuân thủ vào kỹ nghệ yêu cầu từ sớm \citep{kosenkov2025}. Nghiên cứu vì vậy sử dụng chuỗi \emph{nguồn/vị trí $\rightarrow$ ràng buộc $\rightarrow$ quyết định hoặc thành phần kiến trúc $\rightarrow$ tiêu chí kiểm tra}. Đánh giá kiến trúc được thực hiện bằng các kịch bản nhằm làm lộ rủi ro, điểm nhạy và đánh đổi; cách tiếp cận này lấy cảm hứng từ ATAM nhưng không tuyên bố thực hiện đầy đủ quy trình ATAM \citep[pp.~3, 5--7]{gallagher2000}.

\subsection{Liên thông và quản trị dữ liệu như ràng buộc thiết kế}
\label{subsec:vietnam-frameworks}

Khung liên thông châu Âu (EIF) phân biệt bốn lớp pháp lý, tổ chức, ngữ nghĩa và kỹ thuật \citep[Hình~3, p.~22; pp.~27--30]{eif2017}; Quy định (EU) 2024/903 tiếp tục sử dụng EIF như một công cụ hỗ trợ đánh giá liên thông \citep{euact2024}. Trong nghiên cứu này, EIF chỉ là khung phân tích phương pháp. Ở Việt Nam, QĐ 292/QĐ-BKHCN tổ chức kiến trúc theo các mô hình tham chiếu nghiệp vụ, dữ liệu, ứng dụng, công nghệ và an toàn thông tin \citep{qd292_2025}, còn QĐ 2439/QĐ-TTg đặt nền cho quản trị và chuẩn hóa dữ liệu quốc gia \citep{qd2439_2025}.

Các ràng buộc dữ liệu hiện hành làm cho bài toán không thể giải bằng một CSDL nghiệp vụ đơn nhất. NĐ 278/2025/NĐ-CP yêu cầu kết nối/chia sẻ qua hạ tầng điều phối, tuân thủ khung và từ điển dữ liệu, xác định nguồn tin cậy cho dữ liệu chủ và hỗ trợ truy vấn hoặc đồng bộ qua hạ tầng dùng chung \citep{nd278_2025}. QĐ 1973/QĐ-BKHCN đồng thời xác định dữ liệu tổ chức/doanh nghiệp được nhận từ CSDL quốc gia về doanh nghiệp nhưng vẫn liệt kê DN KH\&CN và DN KNST trong dữ liệu chủ chuyên ngành của Bộ \citep{qd1973_2026}. Hệ quả kiến trúc là phải tách \emph{định danh pháp lý} khỏi \emph{trạng thái chuyên ngành} và xác định nguồn có thẩm quyền theo nhóm thuộc tính, thay vì gán một nguồn duy nhất cho toàn bộ thực thể doanh nghiệp.

Định danh, bảo vệ dữ liệu và liên thông cũng tạo các ranh giới riêng. NĐ 69/2024/NĐ-CP và NĐ 169/2025/NĐ-CP yêu cầu sử dụng các nguồn định danh/xác thực có thẩm quyền và đáp ứng điều kiện kết nối \citep{nd69_2024,nd169_2025}; Luật Bảo vệ dữ liệu cá nhân và NĐ 356/2025/NĐ-CP đặt yêu cầu về giới hạn truy cập, biện pháp quản lý/kỹ thuật và vòng đời dữ liệu \citep{luatbvdlcn2025,nd356_2025}. Do đó, SRA cần tách xác thực danh tính khỏi phân quyền nghiệp vụ, đồng thời duy trì nguồn gốc, phiên bản, nhật ký và biên tin cậy cho các luồng dữ liệu liên cơ quan.

\subsection{Hệ thống hiện hữu và ranh giới kế thừa}
\label{subsec:existing-systems}

Định vị SRA đòi hỏi trả lời minh thị nó kế thừa gì từ hạ tầng số đã vận hành hoặc đã được quy định. Trong phạm vi liên quan tới DN KH\&CN và DN KNST, năm hệ thống cần được định vị: Nền tảng số quản lý KH,CN\&ĐMST quốc gia; Hệ thống quản lý nhiệm vụ KH\&CN trực tuyến của Bộ; hệ thống thủ tục hành chính/dịch vụ công; CSDL quốc gia về doanh nghiệp; và các CSDL chuyên ngành theo danh mục QĐ 1762/QĐ-BKHCN. Bảng~\ref{tab:existing-systems} tóm tắt chức năng, dữ liệu thuộc thẩm quyền và quan hệ của từng hệ thống; văn bản dưới đây làm rõ hai điểm bảng không thể hiện hết: quan hệ với nền tảng quốc gia, và trạng thái nguồn của hệ thống quản lý nhiệm vụ KH\&CN.

Luật KH,CN\&ĐMST giao Nhà nước xây dựng, vận hành và duy trì Nền tảng số quản lý KH,CN\&ĐMST quốc gia tập trung, thống nhất, kết nối các chủ thể liên quan \citep[khoản 3 Điều 20]{luatkhcndmst2025}; NĐ 268 cụ thể hóa bằng yêu cầu UBND cấp tỉnh định kỳ cập nhật dữ liệu vòng đời Giấy chứng nhận DN KH\&CN lên nền tảng đó \citep[khoản 3 Điều 48]{nd268_2025}. Nghĩa vụ này thuộc về UBND cấp tỉnh, không phải một hệ thống phần mềm; nền tảng đề xuất -- một \emph{kiến trúc chuyên ngành cấp Bộ} -- chỉ là phương tiện kỹ thuật hỗ trợ cơ quan có thẩm quyền thực hiện nghĩa vụ đó, không phải hệ thống song song hay thay thế Nền tảng số quốc gia. Trong kiến trúc này, nút tỉnh đang có quyết định công nhận/chứng nhận hiệu lực giữ vai trò hệ thống lưu bản ghi nghiệp vụ chính (SoR) cho phạm vi thẩm quyền của mình; lớp trung ương của nền tảng, cũng như Nền tảng số quốc gia, chỉ giữ mô hình đọc, tổng hợp và quan sát liên ngành, không giữ quyền ghi gốc -- nghiệp vụ công nhận/chứng nhận DN KH\&CN và DN KNST mà chưa hệ thống nào khác đảm nhiệm đầy đủ.

Riêng Hệ thống quản lý nhiệm vụ KH\&CN trực tuyến của Bộ -- được phản biện của nghiên cứu trước nêu qua miền vận hành stm.mst.gov.vn -- nay neo được bằng hai căn cứ có văn bản thay vì chỉ mô tả chức năng vận hành công khai. Thứ nhất, nhóm dữ liệu giao dịch ``Khoa học và Công nghệ'' mà QĐ 1973 mô tả -- hồ sơ đề xuất, tuyển chọn, phê duyệt, báo cáo tiến độ/kết quả, nghiệm thu nhiệm vụ \citep{qd1973_2026} -- đúng là phạm vi nghiệp vụ của hệ thống này. Thứ hai, danh mục QĐ 1762 đăng ký CSDL Nhiệm vụ khoa học và công nghệ do Cục Thông tin, Thống kê chủ trì xây dựng và cập nhật \citep[Phụ lục, STT 43]{qd1762_2025}. Với hai căn cứ này, lập luận nền tảng đề xuất không thiết kế lại quy trình quản lý nhiệm vụ KH\&CN đứng vững bằng văn bản; tên miền stm.mst.gov.vn chỉ được nêu như một ví dụ vận hành, không phải căn cứ.

Ba hệ thống còn lại có văn bản dẫn chiếu trực tiếp hơn. Hệ thống TTHC/dịch vụ công giữ kênh tiếp nhận và trả kết quả hồ sơ \citep{nd268_2025,nd278_2025}, đã phân tích riêng ở Mục~\ref{sec:interop}. CSDL quốc gia về doanh nghiệp là nguồn có thẩm quyền cho định danh pháp lý doanh nghiệp \citep{luatdulieu2024,nd278_2025,qd1973_2026}. Danh mục QĐ 1762 liệt kê 59 CSDL chuyên ngành của Bộ KH\&CN, mỗi CSDL có đơn vị chủ trì riêng, trong đó có CSDL Doanh nghiệp khởi nghiệp \citep[Phụ lục, STT 38]{qd1762_2025}. Đáng chú ý, danh mục chỉ đăng ký CSDL Doanh nghiệp khởi nghiệp chứ chưa có mục tương ứng cho DN KH\&CN; đây là điểm chẩn đoán \textsc{cần làm rõ}, không phải xung đột khóa thiết kế, đã được ghi nhận tại Mục~S10.3 và S11 (C07) của tài liệu bổ trợ. Với cả ba hệ thống, nền tảng đề xuất không thiết kế lại: mỗi hệ thống đã có chủ quản và nghĩa vụ báo cáo riêng theo Luật Dữ liệu và QĐ 1973 \citep{luatdulieu2024,qd1973_2026}; tạo bản sao song song sẽ vi phạm nguyên tắc một nguồn có thẩm quyền cho mỗi loại dữ liệu.

Khoảng trống mà SRA thực sự đóng góp nằm ở năm điểm chưa hệ thống nào đảm nhiệm đầy đủ: (i) nghiệp vụ công nhận/chứng nhận DN KH\&CN và DN KNST theo NĐ 268; (ii) hồ sơ năng lực tổng hợp, kết hợp dữ liệu tự khai với bằng chứng đối chiếu từ các nguồn trên; (iii) đối soát, chuẩn hóa dữ liệu khi một doanh nghiệp xuất hiện trên nhiều CSDL với định danh/trạng thái khác nhau; (iv) quản trị thẩm quyền theo nhóm thuộc tính và thời gian thay vì gán tĩnh theo địa bàn hoặc CSDL; và (v) quan hệ hệ sinh thái doanh nghiệp--quỹ--chuyên gia--viện/trường vượt địa giới hành chính \citep[pp.~27--28]{fischer2022}; \citep[pp.~1483--1484]{noak2025}. Nguyên tắc thiết kế theo đó là kế thừa nơi đã có chủ quản rõ ràng, chỉ khác biệt hóa ở khoảng trống thật được xác lập minh thị.

{\vjsttablefont
\renewcommand{\arraystretch}{1.03}
\setlength{\tabcolsep}{3pt}
\begin{longtable}{@{}p{2.6cm}p{4.0cm}p{3.9cm}p{2.9cm}@{}}
\caption{Hệ thống hiện hữu và ranh giới kế thừa của nền tảng đề xuất}
\label{tab:existing-systems}\\
\toprule
\textbf{Hệ thống hiện hữu} & \textbf{Chức năng/dữ liệu đã phủ} & \textbf{Quan hệ với nền tảng đề xuất} & \textbf{Căn cứ} \\
\midrule
\endfirsthead
\multicolumn{4}{c}{\tablename\ \thetable\ -- tiếp theo}\\
\toprule
\textbf{Hệ thống hiện hữu} & \textbf{Chức năng/dữ liệu đã phủ} & \textbf{Quan hệ với nền tảng đề xuất} & \textbf{Căn cứ} \\
\midrule
\endhead
\bottomrule
\endfoot

Nền tảng số quản lý KH,CN\&ĐMST quốc gia & Tổng hợp, kết nối và quan sát liên ngành dữ liệu KH,CN\&ĐMST ở cấp quốc gia; đầu mối nhận cập nhật định kỳ từ các bộ/tỉnh. & UBND cấp tỉnh (SoR nghiệp vụ) dùng nền tảng đề xuất làm phương tiện kỹ thuật để định kỳ cập nhật dữ liệu vòng đời Giấy chứng nhận DN KH\&CN lên nền tảng này; nền tảng đề xuất không giữ quyền ghi cấp Bộ và không phải hệ thống song song. & Luật 93/2025/QH15, khoản 3 Điều 20; NĐ 268, khoản 3 Điều 48 \citep{luatkhcndmst2025,nd268_2025}. \\
\addlinespace
Hệ thống quản lý nhiệm vụ KH\&CN trực tuyến của Bộ (ví dụ vận hành: miền stm.mst.gov.vn) & Đề xuất, tuyển chọn, phê duyệt, triển khai, theo dõi, nghiệm thu nhiệm vụ KH\&CN. & Nền tảng đề xuất tham chiếu kết quả nhiệm vụ KH\&CN như bằng chứng hồ sơ công nhận/chứng nhận; không thiết kế lại quy trình quản lý nhiệm vụ. & Phạm vi dữ liệu giao dịch KH\&CN theo QĐ 1973 \citep{qd1973_2026}; CSDL Nhiệm vụ khoa học và công nghệ, Cục Thông tin, Thống kê chủ trì \citep[Phụ lục, STT 43]{qd1762_2025}. \\
\addlinespace
Hệ thống thủ tục hành chính/dịch vụ công & Tiếp nhận hồ sơ, trả kết quả TTHC cho công dân/doanh nghiệp. & Nền tảng đề xuất trao đổi hồ sơ/kết quả qua định danh dùng chung; không tạo kênh nộp hồ sơ song song. & NĐ 268, khoản 1 Điều 50; NĐ 278 \citep{nd268_2025,nd278_2025}. \\
\addlinespace
CSDL quốc gia về doanh nghiệp & Định danh pháp lý doanh nghiệp. & Nền tảng đề xuất tham chiếu/đồng bộ định danh; không tạo master định danh doanh nghiệp song song. & Luật Dữ liệu; NĐ 278; QĐ 1973 \citep{luatdulieu2024,nd278_2025,qd1973_2026}. \\
\addlinespace
CSDL chuyên ngành theo danh mục QĐ 1762 & 59 CSDL chuyên ngành thuộc phạm vi quản lý của Bộ KH\&CN, mỗi CSDL có đơn vị chủ trì riêng. & Nền tảng đề xuất tham chiếu/xác minh CSDL liên quan khi cần bằng chứng chuyên ngành; không thay thế CSDL đã có chủ quản. & QĐ 1762/QĐ-BKHCN, Phụ lục \citep{qd1762_2025}. \\
\end{longtable}
}

\subsection{Khoảng trống thiết kế và hệ quả đối với SRA}
\label{subsec:design-gap}

Tập nguồn của nghiên cứu được mã hóa thành 10 họ ràng buộc C01--C10 và 128 dòng truy vết nguyên tử; danh mục đầy đủ được chuyển sang tài liệu bổ trợ. Ở mức khái quát, các ràng buộc tạo bốn sức ép kiến trúc đồng thời. Thứ nhất, NĐ 268/2025/NĐ-CP đặt phần lớn quyết định công nhận/chứng nhận ở cấp tỉnh nhưng yêu cầu cập nhật, thông báo và tổng hợp liên cấp, nên quyền ghi nghiệp vụ phải bám thẩm quyền tỉnh trong khi trung ương vẫn cần quan sát nhất quán \citep{nd268_2025}. Thứ hai, Luật Dữ liệu, NĐ 278 và QĐ 1973 yêu cầu tránh tạo lại dữ liệu đã có nguồn chính thức, duy trì chất lượng, nguồn gốc và quản trị chia sẻ \citep{luatdulieu2024,nd278_2025,qd1973_2026}. Thứ ba, NĐ 268 quy định các trạng thái, quyết định và mốc pháp lý chung nhưng việc tổ chức xử lý thuộc cơ quan có thẩm quyền ở địa phương, tạo nhu cầu tách đường cơ sở pháp lý khỏi cấu hình bước nội bộ \citep{nd268_2025}. Thứ tư, nghiên cứu hệ sinh thái khởi nghiệp cho thấy các quan hệ thực tế có thể vượt địa giới hành chính \citep[pp.~27--28]{fischer2022}; \citep[pp.~1483--1484]{noak2025}, vì vậy biên thẩm quyền không được dùng như biên quan hệ dữ liệu.

Các khung kiến trúc và dữ liệu hiện có giúp xác định góc nhìn, nguyên tắc liên thông và yêu cầu quản trị, nhưng không tự quyết định cách tổ chức đồng thời bốn sức ép trên thành một SRA cụ thể. Khoảng trống mà bài giải quyết vì thế không phải là xây dựng một danh mục yêu cầu mới, mà là xác lập: (i) lõi ổn định và điểm biến thiên của kiến trúc tỉnh--trung ương; (ii) mô hình thẩm quyền dữ liệu ở mức nhóm thuộc tính và theo thời gian; (iii) ranh giới trách nhiệm/liên thông độc lập với topology vật lý; và (iv) cách thử thách các quyết định đó bằng kịch bản trước khi triển khai. Đây là cơ sở trực tiếp cho thiết kế nghiên cứu ở Mục~\ref{sec:method} và kiến trúc đề xuất ở Mục~\ref{sec:arch}.
